\documentclass[a4paper,11pt]{book}

\usepackage{polyglossia}
\usepackage{fontspec}
\usepackage{geometry}
\usepackage[intlimits]{amsmath}
\usepackage{graphicx}
\usepackage{indentfirst}
\usepackage{url}

\geometry{verbose,tmargin=4cm,bmargin=3cm,lmargin=3cm,rmargin=2cm,headheight=0.8cm,headsep=1cm,footskip=0.5cm}
\setmainlanguage{czech}
\setmainfont{TeX Gyre Termes}
\setcounter{secnumdepth}{3}

\usepackage{setspace}

\begin{document}
\def\documentdate{7. července 2017}

\pagestyle{empty}
{\centering

\noindent %
\begin{minipage}[c]{3cm}%
\noindent \begin{center}
\includegraphics[width=3cm,height=3cm,keepaspectratio]{images/titlepage/cvut}
\par\end{center}%
\end{minipage}%
\begin{minipage}[c]{0.6\linewidth}%
\begin{center}
\textsc{\large{}České vysoké učení technické v Praze}{\large{}}\\
{\large{}Fakulta jaderná a fyzikálně inženýrská}
\par\end{center}%
\end{minipage}%
\begin{minipage}[c]{3cm}%
\noindent \begin{center}
\includegraphics[width=3cm,height=3cm,keepaspectratio]{images/titlepage/fjfi}
\par\end{center}%
\end{minipage}

\vspace{3cm}


\textbf{\huge{}Název práce česky název práce česky název práce česky
název práce česky}{\huge \par}

\vspace{1cm}


\selectlanguage{american}%
\textbf{\huge{}Title of the Work in English Title of the Work in English
Title of the Work in English}{\huge \par}

\selectlanguage{czech}%
\vspace{2cm}


{\large{}Bakalářská práce}{\large \par}

}

\vfill{}

\begin{list}{}{
	\settowidth{\labelwidth}{MMMMMMMMM}
 \setlength{\leftmargin}{\labelwidth}
 \addtolength{\leftmargin}{\labelsep}
	\renewcommand{\makelabel}[1]{#1\hfil}}
\item [{Autor:}] \textbf{Jméno Autora}
\item [{Vedoucí~práce:}] \textbf{prof. Ing. Jméno Školitele, DrSc.}
\item [{Konzultant:}] \textbf{doc. RNDr. Jméno Konzultanta, CSc. }(pouze
pokud konzultant byl jmenován.)
\item [{Akademický~rok:}] 2016/2017

\end{list}
\newpage{}

~\newpage{}

~

\vfill{}


\begin{center}
- Zadání práce -
\par\end{center}

\vfill{}


~\newpage{}

~

\vfill{}


\begin{center}
- Zadání práce (zadní strana) -
\par\end{center}

\vfill{}


~\newpage{}

\noindent \emph{\Large{}Poděkování:}{\Large \par}

\noindent Chtěl bych zde poděkovat především svému školiteli ...................
za pečlivost, ochotu, vstřícnost a odborné i lidské zázemí při vedení
mé diplomové práce. Dále děkuji svému konzultantovi ................
za ................

\vfill

\noindent \emph{\Large{}Čestné prohlášení:}{\Large \par}

\noindent Prohlašuji, že jsem tuto práci vypracoval samostatně a uvedl
jsem všechnu použitou literaturu.

\bigskip{}


\noindent V Praze dne \documentdate\hfill{}Jméno Autora

\vspace{2cm}


\newpage{}

~\newpage{}

\begin{onehalfspace}
\noindent \emph{Název práce:}

\noindent \textbf{Název práce}
\end{onehalfspace}

\bigskip{}


\noindent \emph{Autor:} Jméno Autora

\bigskip{}


\noindent \emph{Obor:} Celý název oboru (nikoliv zkratka)\bigskip{}


\noindent \emph{Zaměření:} Celý název zaměření (Pokud obor neobsahuje
zaměření, tuto řádku odstranit.)

\bigskip{}


\noindent \emph{Druh práce:} Bakalářská práce

\bigskip{}


\noindent \emph{Vedoucí práce:} prof. Ing. Jméno Školitele, DrSc.,
pracoviště školitele (název instituce, fakulty, katedry...)

\bigskip{}


\noindent \emph{Konzultant:} doc. RNDr. Jméno Konzultanta, CSc., pracoviště
konzultanta. Pouze pokud konzultant byl jmenován.

\bigskip{}


\noindent \emph{Abstrakt:} Abstrakt max. na 10 řádků. Abstrakt max.
na 10 řádků. Abstrakt max. na 10 řádků. Abstrakt max. na 10 řádků.
Abstrakt max. na 10 řádků. Abstrakt max. na 10 řádků. Abstrakt max.
na 10 řádků. Abstrakt max. na 10 řádků. Abstrakt max. na 10 řádků.
Abstrakt max. na 10 řádků. Abstrakt max. na 10 řádků. Abstrakt max.
na 10 řádků. Abstrakt max. na 10 řádků. Abstrakt max. na 10 řádků.
Abstrakt max. na 10 řádků. Abstrakt max. na 10 řádků. Abstrakt max.
na 10 řádků. Abstrakt max. na 10 řádků. Abstrakt max. na 10 řádků.
Abstrakt max. na 10 řádků. Abstrakt max. na 10 řádků. Abstrakt max.
na 10 řádků. Abstrakt max. na 10 řádků. Abstrakt max. na 10 řádků.
Abstrakt max. na 10 řádků. Abstrakt max. na 10 řádků. Abstrakt max.
na 10 řádků. Abstrakt max. na 10 řádků. Abstrakt max. na 10 řádků. 

\bigskip{}


\noindent \emph{Klíčová slova:} klíčová slova (nebo výrazy) seřazená
podle abecedy a oddělená čárkou

\vfill{}
~

\selectlanguage{american}%
\begin{onehalfspace}
\noindent \emph{Title:}

\noindent \textbf{Title of the Work}
\end{onehalfspace}

\bigskip{}


\noindent \emph{Author:} Author's Name

\bigskip{}


\noindent \emph{Abstract:} Max. 10 lines of English abstract text.
Max. 10 lines of English abstract text. Max. 10 lines of English abstract
text. Max. 10 lines of English abstract text. Max. 10 lines of English
abstract text. Max. 10 lines of English abstract text. Max. 10 lines
of English abstract text. Max. 10 lines of English abstract text.
Max. 10 lines of English abstract text. Max. 10 lines of English abstract
text. Max. 10 lines of English abstract text. Max. 10 lines of English
abstract text. Max. 10 lines of English abstract text. Max. 10 lines
of English abstract text. Max. 10 lines of English abstract text.
Max. 10 lines of English abstract text. Max. 10 lines of English abstract
text. Max. 10 lines of English abstract text. Max. 10 lines of English
abstract text. Max. 10 lines of English abstract text. Max. 10 lines
of English abstract text. Max. 10 lines of English abstract text.
Max. 10 lines of English abstract text. Max. 10 lines of English abstract
text. Max. 10 lines of English abstract text.

\bigskip{}


\noindent \emph{Key words:} keywords in alphabetical order separated
by commas

\newpage{}

~\newpage{}

\pagestyle{plain}

\tableofcontents{}

\newpage{}


\chapter*{Úvod}

\addcontentsline{toc}{chapter}{Úvod}

Text úvodu....


\chapter{Název první kapitoly}

\pagestyle{headings}


\chapter*{Závěr}

\pagestyle{plain}

\addcontentsline{toc}{chapter}{Závěr}

Text závěru....
\begin{thebibliography}{1}
\bibitem{Allen-Cahn}S. Allen, J. W. Cahn: \emph{A microscopic theory
for antiphase boundary motion and its application to antiphase domain
coarsening}. Acta Metall., 27:1084-1095, 1979.

\bibitem{CINECA}G. Ballabio et al.: \emph{High Performance Systems
User Guide}. High Performance Systems Department, CINECA, Bologna,
2005. \url{www.cineca.it}

\bibitem{rumpf3}J. Becker, T. Preusser, M. Rumpf: \emph{PDE methods
in flow simulation post processing}. Computing and Visualization in
Science, 3(3):159-167, 2000.\end{thebibliography}

\end{document}
